{V.1.1}					&	% Nomor Revisi
{12 Agustus 2026}			&	% Tanggal Revisi
{\noindent					% Isi Perbaikan Dokumen
    \begin{itemize}[leftmargin=*, nosep]
        \item Menambahkan penjelasan terapan \textit{State-Space} sebagai landasan \textit{Physics Engine} dan penetapan MILP sebagai \textit{Supervisory Discrete-Event Controller} pada BAB 4.
        \item Menambahkan rasionalisasi penggunaan estimasi $\Delta P_{deficit}$ melalui deviasi frekuensi (RoCoF) untuk kendali \textit{transient} pada BAB 4.
        \item Menambahkan sub-bab khusus perbandingan fundamental antara UFLS Konvensional (\textit{blind shedding}) dengan MILP-OLS (\textit{Adaptive UFLS}) pada BAB 5.
        \item Mempertegas penggunaan topologi tunggal (\textit{Copper-Plate}) sebagai \textit{trade-off} performa 10 FPS, dan menetapkan pengujian skala masif (\textit{test system} IEEE) sebagai arah pengembangan lanjutan pada BAB 7.
    \end{itemize}
}